
\section{带标记圆盘模空间及模空间奇异层的定向}
\label{sec:8-3}

本节先赋予带标记点伪全纯圆盘的模空间一个定向；然后，把在一点
粘合的两个圆盘之并上的伪全纯映射模空间，看作（不可约）伪全纯圆盘
模空间的边界，描述它的定向。

首先如下固定 $\operatorname{PSL}(2;\mathbb R)$ 的定向。

\begin{statement}{约定}{8.3.1}\label{conv:8-3-1}
回忆，我们以逆时针方向作为 $\partial D^2$ 的定向约定。取
$\partial D^2$ 上三个互异的标记点 $z_0,z_1,z_2$，其顺序与边界的
逆时针定向相容。通过
\[
 g\longmapsto(g\cdot z_0,g\cdot z_1,g\cdot z_2)
\]
把 $\operatorname{PSL}(2;\mathbb R)$ 嵌入
$\partial D^2\times\partial D^2\times\partial D^2$，并把它作为该乘积
的开子集赋予定向。若固定前两个标记点 $z_0,z_1$，二者的共同稳定子群
为 $\mathbb R$；其定向规定如下：对 $p\ne z_0,z_1$，轨道
$t\mapsto t\cdot p$ 在 $t\to+\infty$ 时收敛到 $z_0$，在
$t\to-\infty$ 时收敛到 $z_1$。为使
$\operatorname{PSL}(2;\mathbb R)$ 在 $\partial D^2$ 上的作用成为右作用，
我们置
\[
 z\cdot g:=g^{-1}z.
\]
因此，对 $\operatorname{PSL}(2;\mathbb R)$ 和 $\mathbb R$，我们采用
上述定向的反向。
\end{statement}

令
\[
 \widetilde{\mathcal M}^{\mathrm{reg}}(\beta)
 =\widetilde{\mathcal M}^{\mathrm{reg}}(L;\beta)
\]
为从单位圆盘出发、代表同伦类
$\beta\in\pi_2(M,L)$ 的伪全纯映射空间。（为简化记号，我们往往省略
$L$。）第 8.3、8.4 节处理
$\widetilde{\mathcal M}^{\mathrm{reg}}(\beta)$ 的全部连通分支；第 8.5 节
则限于主分支，并给出滤过 $A_\infty$ 运算中的符号。回忆，由定理
\ref{thm:8-1-1}，一个相对 spin 结构确定
$\widetilde{\mathcal M}^{\mathrm{reg}}(\beta)$ 的定向。记
\[
 \mathcal M_{m+1}^{\mathrm{reg}}(\beta)
 =\Bigl(\widetilde{\mathcal M}^{\mathrm{reg}}(\beta)
 \times(\partial D^2)^{m+1}\Bigr)
 /\operatorname{PSL}(2;\mathbb R).
\]
把 $+1\in\partial D^2$ 指定为第 $0$ 个标记点 $z_0$，把 $-1$ 指定为
第 $1$ 个标记点 $z_1$，并置%
\translatornote{源书下式右端写成
$\widetilde{\mathcal M}(\beta)\times(\partial D^2)^{m-1}$，漏去上标
$\mathrm{reg}$；下一页对开子集的严格说明、随后同一等式及全文其余记号
均保留该上标，故这里补回。}
\[
 \widetilde{\mathcal M}_{m+1}^{\mathrm{reg}}(\beta)
 =\widetilde{\mathcal M}^{\mathrm{reg}}(\beta)
  \times(\partial D^2)^{m-1}.
\]
这个空间可视为从单位圆盘出发、代表 $\beta$，并在圆盘边界上带有
$m+1$ 个标记点 $(z_0,\ldots,z_m)$ 的伪全纯映射空间，其中固定
$z_0=1$、$z_1=-1$。

严格来说，我们要求所有标记点互异。因此
$\widetilde{\mathcal M}_{m+1}^{\mathrm{reg}}(\beta)$ 并不等于
$\widetilde{\mathcal M}^{\mathrm{reg}}(\beta)
\times(\partial D^2)^{m-1}$ 本身，而是它的一个开子集。不过在讨论定向
问题时，为简便起见，我们常把它写成
\[
 \widetilde{\mathcal M}_{m+1}^{\mathrm{reg}}(\beta)
 =\widetilde{\mathcal M}^{\mathrm{reg}}(\beta)
  \times(\partial D^2)^{m-1}.
\]
于是
$\mathcal M_{m+1}^{\mathrm{reg}}(\beta)$ 与
\[
 \widetilde{\mathcal M}_{m+1}^{\mathrm{reg}}(\beta)
 /\operatorname{Aut}(D^2;z_0,z_1)
\]
相识别；这里 $\operatorname{Aut}(D^2;z_0,z_1)$ 是固定 $z_0,z_1$ 的
$D^2$ 的双全纯自同构群。

由于
$\operatorname{PSL}(2;\mathbb R)/\operatorname{Aut}(D^2;z_0,z_1)$
的维数为偶数，把模空间作为
$\operatorname{PSL}(2;\mathbb R)$ 的商所得的定向，与把它作为
$\operatorname{Aut}(D^2;z_0,z_1)$ 的商所得者相同；比较式
\eqref{eq:8-2-1-2}。记 $\mathbb R_\beta$ 为作用在
$\widetilde{\mathcal M}_{m+1}^{\mathrm{reg}}(\beta)$ 上的自同构群
\[
 \operatorname{Aut}(D^2;z_0=+1,z_1=-1).
\]
（按照我们对从右作用于 $D^2$ 的
$\operatorname{PSL}(2;\mathbb R)$ 所作的定向定义，这个
$\mathbb R$-作用的轨道在 $t\to+\infty$ 时收敛到 $-1$，在
$t\to-\infty$ 时收敛到 $+1$。）注意
\[
 \dim\widetilde{\mathcal M}_{m+1}^{\mathrm{reg}}(\beta)
 =n+\mu(\beta)+m-1.
\]

这些空间的定向，由定向空间的乘积以及对自同构群取商来定义（见第
8.2 节约定~\ref{conv:8-2-1}(2)）。所以在讨论定向时，我们将简写为
\begin{equation}
 \widetilde{\mathcal M}_{m+1}^{\mathrm{reg}}(\beta)
 =\mathcal M_{m+1}^{\mathrm{reg}}(\beta)\times\mathbb R_\beta,
 \tag{8.3.2}\label{eq:8-3-2}
\end{equation}
并且
\[
 \dim\mathcal M_{m+1}^{\mathrm{reg}}(\beta)
 =n+\mu(\beta)+m-2\equiv n+m\pmod2.
\]
回忆，由于假定拉格朗日子流形可定向，马斯洛夫指标为偶数。

用 $\operatorname{ev}_j^\beta$ 表示第 $j$ 个标记点 $z_j$ 处的评价映射
\[
 \operatorname{ev}_j^\beta\colon
 \mathcal M_{m+1}^{\mathrm{reg}}(\beta)\longrightarrow L
 \qquad(j=0,\ldots,m).
\]

\begin{proposition}{8.3.3}\label{prop:8-3-3}
作为带仓西结构的定向空间，有同构
\[
\begin{aligned}
 \partial\mathcal M_{m+1}^{\mathrm{reg}}(\beta'+\beta'')
 ={}&\bigcup
 (-1)^{(m_1-1)(m_2-1)+(n+m_1)}\\
 &\quad\mathcal M_{m_1+1}^{\mathrm{reg}}(\beta')
 {}_{\operatorname{ev}_1^{\beta'}}\!\times_{\operatorname{ev}_0^{\beta''}}
 \mathcal M_{m_2+1}^{\mathrm{reg}}(\beta''),
\end{aligned}
\]
其中 $n=\dim L$；并集遍历
$\beta=\beta'+\beta''$、$m=m_1+m_2-1$。
\end{proposition}

\begin{remark}{8.3.4}\label{rem:8-3-4}
这里单位圆盘边界上标记点的次序规定如下。粘合全纯圆盘时，把
$\beta'$-圆盘的第 $1$ 个标记点 $z_{1,\beta'}$ 与
$\beta''$-圆盘的第 $0$ 个标记点 $z_{0,\beta''}$ 粘合；磨光以后，这两个
标记点消失。粘合后，$\beta'$-圆盘的第 $0$ 个标记点
$z_{0,\beta'}$ 成为 $(\beta'+\beta'')$-圆盘的第 $0$ 个标记点
$z_{0,\beta'+\beta''}$。先把 $\beta''$-圆盘的其余标记点编号为
$1$ 至 $m_2$，再把 $\beta'$-圆盘的其余标记点编号为
$m_2+1$ 至 $m=m_1+m_2-1$。也就是说，
\[
\begin{aligned}
 z_{j,\beta'+\beta''}&=z_{j,\beta''}
 &&(j=1,\ldots,m_2),\\
 z_{m_2+j,\beta'+\beta''}&=z_{j+1,\beta'}
 &&(j=1,\ldots,m_1-1).
\end{aligned}
\]
我们把这一约定写成
\[
\begin{aligned}
 &\partial D^2_{2,\beta'+\beta''}\times\cdots\times
  \partial D^2_{m,\beta'+\beta''}\\
 &\qquad={}
 \partial D^2_{2,\beta''}\times\cdots\times
 \partial D^2_{m_2,\beta''}\times
 \partial D^2_{2,\beta'}\times\cdots\times
 \partial D^2_{m_1,\beta'},
\end{aligned}
\]
或简写为
\[
 (\partial D^2)^{m-1}
 =(\partial D^2)^{m_2-1}\times(\partial D^2)^{m_1-1}.
\]
\end{remark}

\par\medskip\noindent\begin{minipage}{\textwidth}
\centering
\setlength{\unitlength}{0.86mm}
\begin{picture}(142,91)
  \linethickness{0.42pt}
  \put(40,66){\circle{32}}
  \put(101,66){\circle{32}}
  \put(40,66){\makebox(0,0){$\beta''$}}
  \put(101,66){\makebox(0,0){$\beta'$}}
  \put(56,66){\line(0,1){2.4}}
  \put(56,67){\makebox(0,0)[l]{\scriptsize $z_{0,\beta''}$}}
  \put(50,78){\line(1,1){1.8}}
  \put(51.5,80){\makebox(0,0)[l]{\scriptsize $z_{1,\beta''}$}}
  \put(39,82){\line(0,1){2.4}}
  \put(39,86){\makebox(0,0){\scriptsize $z_{2,\beta''}$}}
  \put(47,52){\line(1,-1){1.8}}
  \put(46,48){\makebox(0,0){\scriptsize $z_{m_2,\beta''}$}}
  \put(85,66){\line(0,1){2.4}}
  \put(83,67){\makebox(0,0)[r]{\scriptsize $z_{1,\beta'}$}}
  \put(117,66){\line(0,1){2.4}}
  \put(119,67){\makebox(0,0)[l]{\scriptsize $z_{0,\beta'}$}}
  \put(110,78){\line(1,1){1.8}}
  \put(113,81){\makebox(0,0)[l]{\scriptsize $z_{2,\beta'}$}}
  \put(100,82){\line(0,1){2.4}}
  \put(100,86){\makebox(0,0){\scriptsize $z_{3,\beta'}$}}
  \put(108,52){\line(1,-1){1.8}}
  \put(108,48){\makebox(0,0){\scriptsize $z_{m_1,\beta'}$}}
  \multiput(23,78)(-1.5,-3.1){7}{\circle*{0.55}}
  \multiput(84,78)(-1.5,-3.1){7}{\circle*{0.55}}
  \put(56,22){\circle{30}}
  \put(86,22){\circle{30}}
  \put(56,22){\makebox(0,0){$\beta''$}}
  \put(86,22){\makebox(0,0){$\beta'$}}
  \put(41,27){\line(-1,1){1.7}}
  \put(38,31){\makebox(0,0)[r]{\scriptsize $z_{2,\beta'+\beta''}$}}
  \put(68,34){\line(1,1){1.7}}
  \put(66,38){\makebox(0,0)[r]{\scriptsize $z_{1,\beta'+\beta''}$}}
  \put(101,22){\line(0,1){2.3}}
  \put(104,23){\makebox(0,0)[l]{\scriptsize $z_{0,\beta'+\beta''}$}}
  \put(94,34){\line(1,1){1.7}}
  \put(96,38){\makebox(0,0)[l]{\scriptsize $z_{m_2+1,\beta'+\beta''}$}}
  \put(86,37){\line(0,1){2.3}}
  \put(86,41){\makebox(0,0){\scriptsize $z_{m_2+2,\beta'+\beta''}$}}
  \put(62,9){\line(1,-1){1.7}}
  \put(58,5){\makebox(0,0){\scriptsize $z_{m_2,\beta'+\beta''}$}}
  \put(93,9){\line(1,-1){1.7}}
  \put(95,5){\makebox(0,0){\scriptsize $z_{m,\beta'+\beta''}$}}
  \multiput(39,34)(-1.4,-2.8){7}{\circle*{0.55}}
\end{picture}
\par\smallskip
\renewcommand{\thefigure}{8.3.1}
\refstepcounter{figure}
\textbf{图~\thefigure}\label{fig:8-3-1}
\end{minipage}\par\medskip


\begin{proof}[命题~\ref{prop:8-3-3} 的证明]
先陈述下列引理。

\begin{lemma}{8.3.5}\label{lem:8-3-5}
胶合映射
\[
 \widetilde{\mathcal M}^{\mathrm{reg}}(\beta')
 {}_{\operatorname{ev}_1^{\beta'}}\!\times_{\operatorname{ev}_0^{\beta''}}
 \widetilde{\mathcal M}^{\mathrm{reg}}(\beta'')
 \longrightarrow
 \widetilde{\mathcal M}^{\mathrm{reg}}(\beta'+\beta'')
\]
在仓西结构的意义下保定向。
\end{lemma}

讨论定向问题时，我们将简写为
\[
 \widetilde{\mathcal M}^{\mathrm{reg}}(\beta'+\beta'')
 =\widetilde{\mathcal M}^{\mathrm{reg}}(\beta')
 {}_{\operatorname{ev}_1^{\beta'}}\!\times_{\operatorname{ev}_0^{\beta''}}
 \widetilde{\mathcal M}^{\mathrm{reg}}(\beta'').
\]
引理~\ref{lem:8-3-5} 的证明放在本节末尾。利用该引理，命题的证明如下：
\[
\begin{aligned}
 &\widetilde{\mathcal M}_{m+1}^{\mathrm{reg}}(\beta'+\beta'')\\
 &={}
 \widetilde{\mathcal M}^{\mathrm{reg}}(\beta'+\beta'')
 \times(\partial D^2)^{m-1}\\
 &={}
 \left(
 \widetilde{\mathcal M}^{\mathrm{reg}}(\beta')
 {}_{\operatorname{ev}_1^{\beta'}}\!\times_{\operatorname{ev}_0^{\beta''}}
 \widetilde{\mathcal M}^{\mathrm{reg}}(\beta'')
 \right)\times(\partial D^2)^{m-1}\\
 &=(-1)^{(m_1-1)(m_2-1)}
 \left(\widetilde{\mathcal M}^{\mathrm{reg}}(\beta')
       \times(\partial D^2)^{m_1-1}\right)\\[-2pt]
 &\qquad{}
 {}_{\operatorname{ev}_1^{\beta'}}\!\times_{\operatorname{ev}_0^{\beta''}}
 \left(\widetilde{\mathcal M}^{\mathrm{reg}}(\beta'')
       \times(\partial D^2)^{m_2-1}\right)\\
 &=(-1)^{(m_1-1)(m_2-1)}
 \left(\mathcal M_{m_1+1}^{\mathrm{reg}}(\beta')
       \times\mathbb R_{\beta'}\right)\\[-2pt]
 &\qquad{}
 {}_{\operatorname{ev}_1^{\beta'}}\!\times_{\operatorname{ev}_0^{\beta''}}
 \left(\mathcal M_{m_2+1}^{\mathrm{reg}}(\beta'')
       \times\mathbb R_{\beta''}\right)\\
 &=(-1)^{(m_1-1)(m_2-1)+(n+m_1)}
 \left(\mathbb R_{\beta'}\times
       \mathcal M_{m_1+1}^{\mathrm{reg}}(\beta')\right)\\[-2pt]
 &\qquad{}
 {}_{\operatorname{ev}_1^{\beta'}}\!\times_{\operatorname{ev}_0^{\beta''}}
 \left(\mathcal M_{m_2+1}^{\mathrm{reg}}(\beta'')
       \times\mathbb R_{\beta''}\right)\\
 &=(-1)^{(m_1-1)(m_2-1)+(n+m_1)}
 \mathbb R_{\beta'}\\[-2pt]
 &\qquad{}\times
 \left(
 \mathcal M_{m_1+1}^{\mathrm{reg}}(\beta')
 {}_{\operatorname{ev}_1^{\beta'}}\!\times_{\operatorname{ev}_0^{\beta''}}
 \mathcal M_{m_2+1}^{\mathrm{reg}}(\beta'')
 \right)\times\mathbb R_{\beta''}\\
 &=(-1)^{(m_1-1)(m_2-1)+(n+m_1)}
 \mathbb R_{\mathrm{out}}\\[-2pt]
 &\qquad{}\times
 \left(
 \mathcal M_{m_1+1}^{\mathrm{reg}}(\beta')
 {}_{\operatorname{ev}_1^{\beta'}}\!\times_{\operatorname{ev}_0^{\beta''}}
 \mathcal M_{m_2+1}^{\mathrm{reg}}(\beta'')
 \right)\times\mathbb R_{\beta'+\beta''}.
\end{aligned}
\]
这里说明上面若干等式。为证明第三行的等式，置
\[
 \widetilde{\mathcal M}^{\mathrm{reg}}(\beta')
 =\widetilde{\mathcal M}^{\mathrm{reg}}(\beta')^\circ\times L,
 \qquad
 \widetilde{\mathcal M}^{\mathrm{reg}}(\beta'')
 =L\times{}^\circ\widetilde{\mathcal M}^{\mathrm{reg}}(\beta'').
\]
由于
\[
\begin{aligned}
 \widetilde{\mathcal M}^{\mathrm{reg}}(\beta')
 \times(\partial D^2)^{m_1-1}
 ={}&(-1)^{(m_1-1)n}
 \widetilde{\mathcal M}^{\mathrm{reg}}(\beta')^\circ
 \times(\partial D^2)^{m_1-1}\times L,
\end{aligned}
\]
故有
\[
\begin{aligned}
 &\left(\widetilde{\mathcal M}^{\mathrm{reg}}(\beta')
        \times(\partial D^2)^{m_1-1}\right)\\[-2pt]
 &\qquad{}
 {}_{\operatorname{ev}_1^{\beta'}}\!\times_{\operatorname{ev}_0^{\beta''}}
 \left(\widetilde{\mathcal M}^{\mathrm{reg}}(\beta'')
       \times(\partial D^2)^{m_2-1}\right)\\
 &={}
 (-1)^{(m_1-1)n}
 \widetilde{\mathcal M}^{\mathrm{reg}}(\beta')^\circ
 \times(\partial D^2)^{m_1-1}\times L\\[-2pt]
 &\qquad{}\times{}^\circ
 \widetilde{\mathcal M}^{\mathrm{reg}}(\beta'')
 \times(\partial D^2)^{m_2-1}\\
 &={}
 (-1)^{(m_1-1)n+(m_1-1)(n+m_2-1)}
 \widetilde{\mathcal M}^{\mathrm{reg}}(\beta')^\circ
 \times L\\[-2pt]
 &\qquad{}\times{}^\circ
 \widetilde{\mathcal M}^{\mathrm{reg}}(\beta'')
 \times(\partial D^2)^{m_2-1}\times(\partial D^2)^{m_1-1}\\
 &=(-1)^{(m_1-1)(m_2-1)}
 \left(
 \widetilde{\mathcal M}^{\mathrm{reg}}(\beta')
 {}_{\operatorname{ev}_1^{\beta'}}\!\times_{\operatorname{ev}_0^{\beta''}}
 \widetilde{\mathcal M}^{\mathrm{reg}}(\beta'')
 \right)\times(\partial D^2)^{m-1}.
\end{aligned}
\]

\par\medskip\noindent\begin{minipage}{\textwidth}
\centering
\setlength{\unitlength}{0.76mm}
\begin{picture}(166,139)
  \linethickness{0.4pt}
  \put(83,135){\makebox(0,0)
    {$ (w,z_i)\longmapsto(w\circ\varphi,\varphi^{-1}(z_i)) $}}

  \put(31,105){\circle{25}}
  \put(56,105){\circle{25}}
  \put(31,105){\makebox(0,0){$\beta''$}}
  \put(56,105){\makebox(0,0){$\beta'$}}
  \put(17,105){\makebox(0,0)[r]{\scriptsize $z_{1,\beta''}$}}
  \put(43.5,105){\circle*{0.75}}
  \put(42,109){\makebox(0,0)[r]{\scriptsize $z_{0,\beta''}$}}
  \put(45,109){\makebox(0,0)[l]{\scriptsize $z_{1,\beta'}$}}
  \put(70,105){\makebox(0,0)[l]{\scriptsize $z_{0,\beta'}$}}
  \qbezier(18,117)(31,127)(43,117)
  \qbezier(44,117)(56,127)(69,117)
  \qbezier(18,93)(31,83)(43,93)
  \qbezier(44,93)(56,83)(69,93)
  \put(42,117){\vector(1,-1){2}}
  \put(45,117){\vector(-1,-1){2}}
  \put(42,93){\vector(1,1){2}}
  \put(45,93){\vector(-1,1){2}}
  \put(43.5,84){\makebox(0,0)
    {$\mathbb R_{\mathrm{out}}=\mathbb R_{(1,-1)}$}}
  \put(43.5,77){\makebox(0,0)
    {$\begin{smallmatrix}\beta'&\beta''\\(1,&-1)\end{smallmatrix}$}}

  \put(110,105){\circle{25}}
  \put(135,105){\circle{25}}
  \put(110,105){\makebox(0,0){$\beta''$}}
  \put(135,105){\makebox(0,0){$\beta'$}}
  \put(96,105){\makebox(0,0)[r]{\scriptsize $z_{1,\beta''}$}}
  \put(149,105){\makebox(0,0)[l]{\scriptsize $z_{0,\beta'}$}}
  \qbezier(97,117)(110,127)(122,117)
  \qbezier(123,117)(135,127)(148,117)
  \qbezier(97,93)(110,83)(122,93)
  \qbezier(123,93)(135,83)(148,93)
  \put(98,117){\vector(-1,-1){2}}
  \put(124,117){\vector(-1,-1){2}}
  \put(98,93){\vector(-1,1){2}}
  \put(124,93){\vector(-1,1){2}}
  \put(122.5,84){\makebox(0,0)
    {$\mathbb R_{\beta'+\beta''}=\mathbb R_{(1,1)}$}}
  \put(122.5,77){\makebox(0,0)
    {$\begin{smallmatrix}\beta'&\beta''\\(1,&1)\end{smallmatrix}$}}
  \put(83,68){\makebox(0,0){稳定映射上的拉回 $\mathbb R$-作用}}

  \put(31,42){\circle{25}}
  \put(56,42){\circle{25}}
  \put(31,42){\makebox(0,0){$\beta''$}}
  \put(56,42){\makebox(0,0){$\beta'$}}
  \put(17,42){\makebox(0,0)[r]{\scriptsize $z_{1,\beta''}$}}
  \put(43.5,42){\circle*{0.75}}
  \put(42,46){\makebox(0,0)[r]{\scriptsize $z_{0,\beta''}$}}
  \put(45,46){\makebox(0,0)[l]{\scriptsize $z_{1,\beta'}$}}
  \put(70,42){\makebox(0,0)[l]{\scriptsize $z_{0,\beta'}$}}
  \qbezier(18,54)(31,64)(43,54)
  \qbezier(44,54)(56,64)(69,54)
  \qbezier(18,30)(31,20)(43,30)
  \qbezier(44,30)(56,20)(69,30)
  \put(19,54){\vector(-1,-1){2}}
  \put(68,54){\vector(1,-1){2}}
  \put(19,30){\vector(-1,1){2}}
  \put(68,30){\vector(1,1){2}}
  \put(43.5,21){\makebox(0,0)
    {$\mathbb R_{\mathrm{out}}=\mathbb R_{(1,-1)}$}}

  \put(110,42){\circle{25}}
  \put(135,42){\circle{25}}
  \put(110,42){\makebox(0,0){$\beta''$}}
  \put(135,42){\makebox(0,0){$\beta'$}}
  \put(96,42){\makebox(0,0)[r]{\scriptsize $z_{1,\beta''}$}}
  \put(122.5,42){\circle*{0.75}}
  \put(121,46){\makebox(0,0)[r]{\scriptsize $z_{0,\beta''}$}}
  \put(124,46){\makebox(0,0)[l]{\scriptsize $z_{1,\beta'}$}}
  \put(149,42){\makebox(0,0)[l]{\scriptsize $z_{0,\beta'}$}}
  \qbezier(97,54)(110,64)(122,54)
  \qbezier(123,54)(135,64)(148,54)
  \qbezier(97,30)(110,20)(122,30)
  \qbezier(123,30)(135,20)(148,30)
  \put(121,54){\vector(1,-1){2}}
  \put(147,54){\vector(1,-1){2}}
  \put(121,30){\vector(1,1){2}}
  \put(147,30){\vector(1,1){2}}
  \put(122.5,21){\makebox(0,0)
    {$\mathbb R_{\beta'+\beta''}=\mathbb R_{(1,1)}$}}
  \put(83,12){\makebox(0,0){标记点上的 $\mathbb R$-作用}}
  \put(43.5,5){\makebox(0,0){$(1,-1)\quad\mathbb R_{\mathrm{out}}$（胶合）}}
  \put(122.5,5){\makebox(0,0){$(1,1)\quad\mathbb R_{\beta'+\beta''}$}}
\end{picture}
\par\smallskip
\renewcommand{\thefigure}{8.3.2}
\refstepcounter{figure}
\textbf{图~\thefigure}\label{fig:8-3-2}
\end{minipage}\par\medskip


在最后一个等式中，$\mathbb R_{\mathrm{out}}$ 和
$\mathbb R_{\beta'+\beta''}$ 分别是
$\mathbb R_{\beta'}\times\mathbb R_{\beta''}$ 中带定向的反对角线与
对角线子流形；也就是说，$\mathbb R_{\mathrm{out}}$ 由
$(1,-1)\in\mathbb R_{\beta'}\times\mathbb R_{\beta''}$ 生成，而
$\mathbb R_{\beta'+\beta''}$ 由 $(1,1)$ 生成。因此
$\mathbb R_{\beta'}\times\mathbb R_{\beta''}$ 的定向，与
$\mathbb R_{\mathrm{out}}\times\mathbb R_{\beta'+\beta''}$ 的定向相同。
因子 $\mathbb R_{\mathrm{out}}$ 可视为全纯圆盘的胶合参数空间，而
$\mathbb R_{\beta'+\beta''}$ 可视为
$\widetilde{\mathcal M}_{m+1}^{\mathrm{reg}}(\beta'+\beta'')$ 上的
$\mathbb R$-作用（见图~\ref{fig:8-3-2}）。由边界定向与商定向的定义，
便得到命题~\ref{prop:8-3-3}。
\end{proof}

\begin{proof}[引理~\ref{lem:8-3-5} 的证明]
设
$u\in\widetilde{\mathcal M}^{\mathrm{reg}}(\beta')$、
$v\in\widetilde{\mathcal M}^{\mathrm{reg}}(\beta'')$ 是满足
$u(-1)=v(1)$ 的全纯圆盘。固定充分大的 $R>0$；利用单位分解和隐函数
定理，分别在 $u$ 的 $-1$ 附近与 $v$ 的 $1$ 附近取半径 $1/R$ 的邻域
进行胶合（见第~\ref{subsec:7-1-3} 小节）。把所得全纯圆盘记为
\[
 u\#_Rv\in
 \widetilde{\mathcal M}^{\mathrm{reg}}(\beta'+\beta'').
\]
$L$ 上的相对 spin 结构给出
$(u|_{\partial D^2})^*TL$、$(v|_{\partial D^2})^*TL$ 和
$((u\#_Rv)|_{\partial D^2})^*TL$ 的稳定平凡化，而且这些平凡化在胶合
下相容。因此引理~\ref{lem:8-3-5} 化归到线性化问题；也就是说，它由
下面的引理~\ref{lem:8-3-10} 推出。
\end{proof}

为陈述引理~\ref{lem:8-3-10}，先准备记号。设 $(E_1,F_1)$、
$(E_2,F_2)$ 是单位圆盘上的复向量丛，其边界上带有全实子丛（亦即定义
\ref{def:2-1-9} 意义下的复丛对）。有多尔博算子
\[
 \bar\partial_{E_i,F_i}\colon
 W^{1,p}(D^2,\partial D^2;E_i,F_i)
 \longrightarrow L^p(D^2;E_i)
 \qquad(i=1,2),
\]
它们都是弗雷德霍姆映射。对 $p>2$ 和 $z\in\partial D^2$，定义评价映射
\[
 \operatorname{ev}_z\colon
 W^{1,p}(D^2,\partial D^2;E_i,F_i)
 \longrightarrow F_i|_z,
 \qquad \operatorname{ev}_z(\zeta)=\zeta(z).
\]
假定给定识别 $F_1|_{-1}\cong F_2|_1$。用它胶合丛对
$(E_1,F_1)$、$(E_2,F_2)$，得到 $(D^2,\partial D^2)$ 上的复丛对
$(E,F)$。（这里的 $D^2$ 由两个 $D^2$ 的副本 $D_1^2,D_2^2$ 在
$-1\in D_1^2$ 与 $+1\in D_2^2$ 处胶合而成。）

如下定义指标的纤维积
\[
 \operatorname{Index}(\bar\partial_{E_1,F_1})
 {}_{\operatorname{ev}_{-1}}\!\times_{\operatorname{ev}_1}
 \operatorname{Index}(\bar\partial_{E_2,F_2}).
\]
取有限维复线性子空间
$\mathcal E_i\subset L^p(D^2;E_i)$，满足
\begin{align}
 \operatorname{Im}(\bar\partial_{E_i,F_i})+\mathcal E_i
 &=L^p(D^2;E_i)\qquad(i=1,2),
 \tag{8.3.6}\label{eq:8-3-6}\\
 \operatorname{ev}_{-1}\colon
 \bar\partial_{E_1,F_1}^{-1}(\mathcal E_1)
 &\longrightarrow F_1|_{-1}\quad\text{为满射},
 \tag{8.3.7.1}\label{eq:8-3-7-1}\\
 \operatorname{ev}_{1}\colon
 \bar\partial_{E_2,F_2}^{-1}(\mathcal E_2)
 &\longrightarrow F_2|_{1}\quad\text{为满射}.
 \tag{8.3.7.2}\label{eq:8-3-7-2}
\end{align}

\begin{definition}{8.3.8}\label{def:8-3-8}
置
\[
\begin{aligned}
 &\bar\partial_{E_1,F_1}^{-1}(\mathcal E_1)
 {}_{\operatorname{ev}_{-1}}\!\times_{\operatorname{ev}_1}
 \bar\partial_{E_2,F_2}^{-1}(\mathcal E_2)\\
 &\quad=\bigl\{(v_1,v_2)\in
 W^{1,p}(D^2,\partial D^2;E_1,F_1)
 \oplus W^{1,p}(D^2,\partial D^2;E_2,F_2)\ \bigm|\\
 &\hspace{8em}
 \bar\partial_{E_i,F_i}(v_i)\in\mathcal E_i\ (i=1,2),\quad
 \operatorname{ev}_{-1}(v_1)=\operatorname{ev}_1(v_2)\bigr\},
\end{aligned}
\]
并定义
\begin{equation}
\begin{aligned}
 &\operatorname{Index}(\bar\partial_{E_1,F_1})
 {}_{\operatorname{ev}_{-1}}\!\times_{\operatorname{ev}_1}
 \operatorname{Index}(\bar\partial_{E_2,F_2})\\
 &\quad=
 \bar\partial_{E_1,F_1}^{-1}(\mathcal E_1)
 {}_{\operatorname{ev}_{-1}}\!\times_{\operatorname{ev}_1}
 \bar\partial_{E_2,F_2}^{-1}(\mathcal E_2)
 -\mathcal E_1-\mathcal E_2.
\end{aligned}
\tag{8.3.9}\label{eq:8-3-9}
\end{equation}
这里的式~\eqref{eq:8-3-9} 是定向虚向量空间的等式。容易看出，其右端
与 $\mathcal E_i$ 的选择无关。%
\translatornote{源书紧前明定在 $-1\in D_1^2$ 与 $+1\in D_2^2$ 胶合，
并把指标纤维积及引理 8.3.10 写成第一因子的
$\operatorname{ev}_{-1}$ 与第二因子的 $\operatorname{ev}_1$；但在
(8.3.7.1)--(8.3.7.2)、本定义及 (8.3.9) 右端又把两个端点对调。
这里依胶合点、式左端和引理三重一致的次序统一恢复。}
\end{definition}

\begin{lemma}{8.3.10}\label{lem:8-3-10}
存在定向虚向量空间的同构
\[
 \operatorname{Index}(\bar\partial_{E,F})
 =\operatorname{Index}(\bar\partial_{E_1,F_1})
 {}_{\operatorname{ev}_{-1}}\!\times_{\operatorname{ev}_1}
 \operatorname{Index}(\bar\partial_{E_2,F_2}).
\]
\end{lemma}

\begin{proof}
先考虑 $(E_1,F_1)$、$(E_2,F_2)$ 都平凡的情形，也就是说，$F_1$、
$F_2$ 分别延拓为圆盘上 $E_1$、$E_2$ 的全实子丛。这时
$\operatorname{Index}(\bar\partial_{E,F})$ 可同
$*\in\partial D^2$ 处的纤维 $F|_*$ 识别。此外可以取
$\mathcal E_i=0$；此时式~\eqref{eq:8-3-7-1}、
\eqref{eq:8-3-7-2} 是保定向同构。因此引理在这一情形成立。

在一般情形中，仿照命题~\ref{prop:8-1-4} 的证明，利用边界上全实子丛
的（稳定）平凡化，把丛压缩到圆盘与黎曼球面的单点并上。注意，被识别
的点是圆盘与黎曼球面的内部点。因此纤维积是在复向量丛的一个纤维上
取得的，而该纤维带有复定向（见图~\ref{fig:8-3-3}）。

\par\medskip\noindent\begin{minipage}{\textwidth}
\centering
\setlength{\unitlength}{0.92mm}
\begin{picture}(130,55)
  \linethickness{0.48pt}
  \qbezier(34,28)(20,43)(7,34)
  \qbezier(7,34)(1,28)(7,20)
  \qbezier(7,20)(20,10)(34,28)
  \qbezier(34,28)(48,43)(60,34)
  \qbezier(60,34)(66,28)(60,20)
  \qbezier(60,20)(48,10)(34,28)
  \put(17,30){\makebox(0,0){$E_1$}}
  \put(51,30){\makebox(0,0){$E_2$}}
  \put(13,46){\makebox(0,0){$F_1$}}
  \put(56,46){\makebox(0,0){$F_2$}}
  \put(68,28){\vector(1,0){9}}
  \qbezier(105,28)(96,43)(88,35)
  \qbezier(88,35)(82,28)(88,20)
  \qbezier(88,20)(96,12)(105,28)
  \qbezier(105,28)(114,43)(123,35)
  \qbezier(123,35)(129,28)(123,20)
  \qbezier(123,20)(114,12)(105,28)
  \qbezier(94,35)(88,54)(97,52)
  \qbezier(97,52)(104,48)(96,31)
  \qbezier(115,35)(109,54)(118,52)
  \qbezier(118,52)(125,48)(117,31)
\end{picture}
\par\smallskip
\renewcommand{\thefigure}{8.3.3}
\refstepcounter{figure}
\textbf{图~\thefigure}\label{fig:8-3-3}
\end{minipage}\par\medskip


为了核对
\[
 \bar\partial_{E,F}^{-1}(\mathcal E_1\oplus\mathcal E_2)
\]
与
\[
 \bar\partial_{E_1,F_1}^{-1}(\mathcal E_1)
 {}_{\operatorname{ev}_{-1}}\!\times_{\operatorname{ev}_1}
 \bar\partial_{E_2,F_2}^{-1}(\mathcal E_2)
\]
上的定向相容，只需处理下列线性胶合问题：带两个内部二重点的亏格
$0$ 有边黎曼曲面（即在圆盘的两个内部点各附着一个
$\mathbb{CP}^1$）上多尔博算子的实指标，同胶合后的黎曼球面上、取胶合
向量丛所得多尔博算子的实指标同构。由于二者都带有来自复结构的自然
定向，它们显然相容。
\end{proof}
